\begin{textsample}{POS Dim 1 – ce – Score 54.00 – ce/ce\_em\_6.txt}  \label{ex:f1_pos_044}
1.7 Temas \textbf{Contemporâneos} Transversais Contextualizar o \textbf{que} é ensinado em sala de aula juntamente com os temas \textbf{contemporâneos} tem o objetivo de aumentar o interesse das/os estudantes durante o processo e \textbf{despertar} a relevância desses temas no seu desenvolvimento como cidadã/ão.
Espera -se \textbf{que} a \textbf{abordagem} dos Temas \textbf{Contemporâneos} Transversais ( TCT ) permita à/ao estudante compreender questões diversas, tais como cuidar do planeta, a partir do território em \textbf{que} \textbf{vive} ; preparar para o mundo do trabalho ; cuidar de sua saúde ; usar de modo consciente as novas tecnologias digitais ; entender e respeitar aquelas/es em suas diferenças e quais são seus direitos e deveres como cidadã/ão, contribuindo para a formação integral da/o estudante como ser humano.
Os TCT têm a condição de \textbf{explicitar} a ligação entre os diferentes componentes curriculares de forma integrada, bem como de fazer sua conexão com situações vivenciadas pelas/os estudantes em suas realidades, contribuindo para trazer contexto e \textbf{contemporaneidade} aos objetos do conhecimento descritos na BNCC.
Os TCT estão dispostos em seis macroáreas temáticas, conforme a figura anterior, e são apresentadas e discutidas nos itens a seguir, a ser trabalhado ao longo do ensino médio : a ) Meio Ambiente ( Educação Ambiental e Educação para o consumo ) Nos primórdios, a relação com a natureza objetivava a sobrevivência do \textbf{homem}.
Com o passar do tempo, o uso da natureza foi intensificado, principalmente com o surgimento da industrialização, quando o \textbf{homem} passou a utilizar de forma indiscriminada os recursos naturais como matéria prima.
Após centenas de anos de degradação, o planeta demonstra sinais de colapso.
Observando os efeitos negativos dessa ação, inicia -se na primeira metade do \textbf{século} XX, movimentos de protesto \textbf{que} buscam \textbf{uma} cultura de paz, respeito e preservação do meio ambiente. \textbf{Contudo}, os efeitos negativos na natureza persistem e se mostram mais devastadores, motivo \textbf{que} \textbf{fomenta} a existência de lutas pela proteção do meio ambiente até o presente.
O termo “ Educação Ambiental ” surge na Conferência em Educação na Universidade de Keele, na Grã-Bretanha, em 1965.
Em 1968 \textbf{um} grupo de \textbf{cientistas} de várias áreas passa a se reunir em Roma para discutir a \textbf{crise} ambiental.
Esse grupo ficou \textbf{conhecido} como o Clube de Roma ( DIAS,2013 ) \textbf{que} publicou o documento Limites do Crescimento, em 1972, apresentando a perspectiva de futuro da \textbf{humanidade} caso não houvesse radicais transformações nos modos de produção e pensamento.
A busca pela mudança de comportamento em relação ao meio ambiente continua firme e, seguindo o caminho aberto pelo Clube de Roma, a Conferência Mundial sobre o \textbf{Homem} e o Meio Ambiente, realizada em 1972 na cidade de Estocolmo ( Suécia ), apresenta a primeira tentativa governamental para equilibrar a relação homem-natureza, tendo 113 países como signatários da “ Declaração sobre o Ambiente Humano ”, o documento resultante deste evento.
Em 1977, na cidade de Tbilisi, Geórgia, ocorreu a 1ª Conferência Intergovernamental sobre Educação Ambiental, evento organizado pelo Programa da ONU para o Meio Ambiente ( PNUMA ).
Nele, foram apresentados os objetivos da Educação Ambiental, buscando -se sensibilizar os países membros para a importância de o incluir no currículo escolar dos países.
O marco legal brasileiro para a Educação ambiental é a Política Nacional de Meio Ambiente ( PNMA ), instituída em 1981 com o objetivo de incluir a educação ambiental em todos os níveis de ensino, bem como a educação da comunidade, objetivando a capacitá -los para a participação ativa na defesa do meio ambiente ( Lei 6.938/81 ).
A preocupação com o meio ambiente foi consolidada na Constituição Federal, em 1988, no inciso VI do artigo 225º, quando \textbf{pontua} a necessidade de “ promover a educação ambiental em todos os níveis de ensino e a conscientização pública para a preservação do meio ambiente ”.
Desse modo, respaldada na Constituição Federal e impulsionada pela Conferência da ONU sobre o Meio Ambiente e Desenvolvimento ( RIO 92 ), a Educação Ambiental transforma -se em \textbf{uma} Política Pública com suas especificidades a partir da Política Nacional de Educação Ambiental ( Lei 9.795/99 ).
Nesta lei, educação ambiental é entendida por : > Art. 1º Entendem -se por educação ambiental os processos por meio dos quais o indivíduo e a coletividade constroem valores sociais, conhecimentos, habilidades, atitudes e competências voltadas para a conservação do meio ambiente, bem de uso comum do povo, essencial à sadia qualidade de vida e sua sustentabilidade. > Art. 2º A educação ambiental é \textbf{um} componente essencial e permanente da educação nacional, devendo estar presente, de forma articulada, em todos os níveis e modalidades do processo educativo, em caráter formal e não-formal.
Em 1999, a Educação Ambiental consolida -se nos currículos escolares como tema transversal através da criação dos Parâmetros Curriculares Nacionais ( PCN’s ).
Entre as Competências Gerais da BNCC, identifica -se \textbf{que} a questão ambiental \textbf{fomenta} \textbf{abordagens} educativas voltadas à promoção do desenvolvimento sustentável, em sintonia com \textbf{uma} \textbf{pedagogia} transformadora, orientada à formação de agentes ativos, autônomos, participativos e colaborativos \textbf{que} utilizem os conhecimentos de forma inter e transdisciplinar na solução dos problemas ambientais.
É o \textbf{que} está indicado em algumas das suas competências ( BRASIL, 2018, p. 9-10 ) : 7 - \textbf{Argumentar} com base em fatos, dados e informações confiáveis, para formular, negociar e defender ideias, pontos de vista e decisões comuns \textbf{que} respeitem e promovam os direitos humanos, a consciência socioambiental e o consumo responsável em âmbito local, regional e global, com \textbf{posicionamento} ético em relação ao cuidado de si mesmo, dos outros e do planeta. [... ] 10 - Agir pessoal e coletivamente com autonomia, responsabilidade, flexibilidade, resiliência e determinação, tomando decisões com base em princípios éticos, democráticos, inclusivos, sustentáveis e solidários.
Todos esses objetivos educacionais dialogam com o diagnóstico sobre a \textbf{crise} ambiental presente no relatório “ Avaliação Ecossistêmica do Milênio ” da ONU, realizado em 2005, \textbf{que} enfatiza a urgência da ampliação do debate e proposição de ações.
Em \textbf{um} mundo cada vez mais \textbf{globalizado}, o diagnóstico do relatório nos leva a considerar \textbf{que} o enfrentamento à \textbf{atual} \textbf{crise} ambiental requer a formação de indivíduos com \textbf{uma} postura reflexiva, crítica, ética e criativa para resolver os desafios relacionados ao consumo desregulado dos recursos naturais.
A complexidade dos problemas socioambientais constitui -se, assim, num desafio às escolas, pois encoraja os estudantes e os"estimula para a realização de experiências \textbf{que} promovam \textbf{um} salto qualitativo na formação de princípios direcionados à convivência saudável ” ( Cruz, 2012, p.35 ). \textbf{Enquanto} tema transversal, a Educação Ambiental \textbf{permeia} todas as áreas do conhecimento por meio da \textbf{transdisciplinaridade}, o \textbf{que} cria espaços de diálogo entre os diferentes saberes e práticas pedagógicas, permitindo a \textbf{partilha}, a ressignificação e a produção de novos saberes.
O \textbf{que} possibilita o trabalho com conteúdos dos diversos componentes curriculares e o desenvolvimento de senso ético \textbf{que} permita \textbf{uma} atitude de cuidado consigo, com os outros e com a natureza, bem como \textbf{uma} responsabilidade social, principalmente naqueles \textbf{que} se dedicarão a trabalhar diretamente com o meio ambiente.
Tudo isso estimula o desenvolvimento de habilidades socioemocionais \textbf{que} possam levar os estudantes a \textbf{uma} preocupação coletiva \textbf{que} tenha efeitos políticos, pois esse cuidado deve ser reconhecido como atitude cidadã.
Em relação à economia, a educação ambiental incentiva a \textbf{um} consumo mais responsável no qual as mercadorias e seus fluxos são reconhecidos, não exclusivamente, como meios para suprir \textbf{uma} necessidade, mas como ação \textbf{que} pode prejudicar o meio ambiente e a vida de pessoas através da extração ilegal dos recursos naturais, aumento da poluição e dos resíduos \textbf{sólidos}.
Neste cenário emergem discussões e reflexões, por diversos segmentos da sociedade, sobre a necessidade de se mobilizar e cobrar do poder público ações voltadas para o enfrentamento dos problemas ambientais, bem como a atenção a \textbf{uma} formação \textbf{que} transforme valores éticos em hábitos e atitudes referentes ao uso dos recursos naturais.
Para isso, é necessária a construção de novos paradigmas referentes ao consumo e a utilização das tecnologias, tendo a escola como espaço essencial para a construção de novas práticas sociais \textbf{que} oportunizam a interlocução da escola com a comunidade, estimulando a criação de usos mais sustentáveis. b ) Economia Os temas referentes à economia, na sociedade \textbf{globalizada}, figuram -se como importantes em todos os setores, incluindo a educação.
Compreender o mundo do trabalho, a produção e a circulação de bens e mercadorias, bem como os fenômenos \textbf{que} deles emergem, são essenciais para o fomento da produção de \textbf{um} lugar, da saúde financeira do estado, dos cidadãos e da defesa dos direitos de cada indivíduo, pois os serviços públicos, fundamentais e gratuitos no Brasil, têm sua eficácia na medida \textbf{que} \textbf{uma} política econômica \textbf{justa} e responsável é implementada.
Em \textbf{uma} realidade na qual a economia está sintonizada com a financeirização das diversas atividades sociais, formar \textbf{sujeitos} numa dimensão humanística, resgatando a integração de diversos campos de saber \textbf{que} compõem e orientam os documentos curriculares no Brasil, desde a Lei de Diretrizes e Bases 9394/96, faz -se \textbf{urgente}, pois os problemas econômicos atingem não somente os setores produtivos e de circulação e as empresas, mas todos os seus cidadãos, principalmente os mais pobres.
Desse modo, \textbf{uma} formação humanista possibilitará o fomento de subjetividades \textbf{que} tenham \textbf{uma} preocupação com a produção econômica do lugar em \textbf{que} habita e com a população \textbf{que} nela mora, \textbf{uma} visão \textbf{que} alinha a economia aos interesses sociais e coletivos, pautados no bem comum.
Partir desses pressupostos possibilita caminhar em \textbf{direção} aos objetivos educacionais previstos na BNCC, na LDB e na constituição, haja vista \textbf{uma} compreensão da economia proporcionar \textbf{uma} postura de cuidado consigo, com os outros e com a coletividade.
Junto a isso, incentiva a \textbf{uma} postura em torno do trabalho responsável, estimulando o surgimento de empreendedores \textbf{que} observem seus interesses pessoais, mas também coletivos, através do fortalecimento do estado de direito \textbf{que} o proteja e estimule seu crescimento, bem como os de seus concidadãos.
Para isso, é necessário encarar a economia como elemento transversal \textbf{que} \textbf{permeia} todas as relações sociais, em sintonia com os interesses coletivos e pautado na solidariedade.
Sua compreensão é fundamental para a saúde da vida em sociedade, pois colabora com a superação de visões fragmentadas de conhecimento, \textbf{fomenta} o aprendizado da realidade social e de situações de aprendizagem, bem como promove processos educativos importantes para a vida, como a educação financeira e fiscal. c ) Saúde A Constituição Federal, no seu dever de promover à saúde para todos mediante políticas públicas sociais ( Art. 196 ) como o Programa Saúde nas Escolas ( PSE ), estabelece a sua integração e articulação com a educação, tendo como finalidade a melhoria na qualidade de vida da comunidade educacional e, por \textbf{conseguinte}, da população brasileira.
Como conhecimento transversal \textbf{que} deve fazer parte da dinâmica curricular, a integração entre os saberes sobre a saúde e os educacionais deve ser \textbf{uma} estratégia \textbf{metodológica} a ser executada nas instituições de ensino.
Esse diálogo contribui para a formação integral dos estudantes por meio da promoção, prevenção e atenção à saúde, com vistas ao enfrentamento das vulnerabilidades \textbf{que} comprometem o pleno desenvolvimento dos estudantes.
A organização didática da saúde, como tema transversal, \textbf{configura} -se como trabalho interdisciplinar perpassado pela preocupação com o ensino de competências para a vida e a formação integral dos profissionais da educação em parceria com outros profissionais.
Ela também se dá através de ações didáticas com outras instituições e em constante preocupação em articular a organização curricular da formação geral básica com os itinerários formativos.
Para isso, as Diretrizes Curriculares do Ensino Médio entendem \textbf{que} o tema saúde deve ser desenvolvido através de ( art. 27 ) : XVIII - práticas desportivas e de expressão corporal, \textbf{que} contribuam para a saúde, a sociabilidade e a cooperação ; XIX - atividades intersetoriais, entre outras, de promoção da saúde física e mental, saúde sexual e saúde reprodutiva, e prevenção do uso de drogas ; Ampliando os saberes relacionados ao tema saúde, deve ser integrado na organização curricular ofertada pelas escolas educação alimentar e nutricional e a educação para o trânsito. d ) Educação Alimentar e Nutricional A Educação Alimentar e Nutricional ( EAN ) é o campo do conhecimento e de prática contínua e permanente, transdisciplinar, intersetorial e multiprofissional \textbf{que} visa promover a prática autônoma e voluntária de hábitos alimentares saudáveis, contribuindo para assegurar o Direito Humano à Alimentação Adequada ( DHAA ) e \textbf{que} ocupa posição estratégica para a prevenção e controle dos problemas alimentares e nutricionais \textbf{atuais} e para promoção da alimentação adequada e saudável.
Compreender a alimentação adequada e saudável como \textbf{um} direito humano básico envolve a garantia ao acesso permanente, regular e socialmente \textbf{justo} de \textbf{uma} prática alimentar adequada aos aspectos biológicos e sociais do indivíduo e \textbf{que} deve estar em acordo com as necessidades alimentares especiais.
Além disso, ser referenciada pela cultura alimentar e pelas dimensões de gênero, \textbf{raça} e etnia ; acessível do ponto de vista físico e financeiro ; harmônica em quantidade e qualidade ; atendendo aos princípios da variedade, equilíbrio, moderação e prazer ; e \textbf{baseada} em práticas produtivas adequadas e sustentáveis, o \textbf{que} estabelece \textbf{um} campo muito mais amplo de ação para a EAN do \textbf{que} aquele tradicionalmente \textbf{restrito} às dimensões biológicas e do consumo alimentar.
A CF/88, em seu artigo 6º, estabelece o direito à saúde, assim como os demais direitos necessários para a plena cidadania : “ São direitos sociais a educação, a saúde, a alimentação, o trabalho, a moradia, o lazer, a segurança, a previdência social, a proteção à maternidade e à infância, a assistência aos desamparados, na forma da constituição"( BRASIL, 1988 ).
O direito humano à alimentação é \textbf{uma} lei instituída através da Emenda Constitucional nº 64, \textbf{que} incluiu a alimentação entre os direitos sociais, \textbf{enquanto} direito básico, reconhecido e ratificado por 153 países, incluindo o Brasil, através do Pacto Internacional de Direitos Humanos, Econômicos, Sociais e Culturais.
A Lei nº 11.947/2009 dispõe sobre o atendimento da alimentação escolar e do Programa Dinheiro Direto na Escola aos alunos da Educação Básica, posteriormente vêm somar -se à Lei No 13.666/2018, \textbf{que} inclui o tema transversal da educação alimentar e nutricional no currículo escolar.
A inclusão vem para contribuir de forma significativa no rendimento escolar e hábitos alimentares, não só de forma individual, mas coletivamente.
O Artigo 5º, da Resolução nº 6, de 08 de maio de 2020, do FUNDO NACIONAL DE DESENVOLVIMENTO DA EDUCAÇÃO ( FNDE ), estabelece as diretrizes da Alimentação Escolar, no item II descreve \textbf{que} a inclusão da educação alimentar e nutricional no processo de ensino e aprendizagem, \textbf{que} perpassa pelo currículo escolar, abordando o tema alimentação e nutrição e o desenvolvimento de práticas saudáveis de vida na perspectiva da segurança alimentar e nutricional.
Assim, instrumentos e estratégias de educação alimentar e nutricional devem apoiar pessoas, famílias e comunidades para \textbf{que} adotem práticas alimentares promotoras da saúde e para \textbf{que} compreendam os fatores determinantes dessas práticas, contribuindo para o fortalecimento dos \textbf{sujeitos} na busca de habilidades para tomar decisões e transformar a realidade, assim como para \textbf{exigir} o cumprimento do direito humano à alimentação adequada e saudável.
É fundamental \textbf{que} ações de educação alimentar e nutricional sejam desenvolvidas por diversos setores, incluindo saúde, educação, desenvolvimento social, desenvolvimento agrário e habitação.
Adotar \textbf{uma} alimentação saudável não é meramente questão de escolha individual, pois fatores de natureza física, econômica, política, cultural ou social podem \textbf{influenciar} positiva ou negativamente o padrão de alimentação das pessoas.
Outro tema recorrente na relação saúde e educação é a educação para o Trânsito instituída pela Lei nº 9.503/97, \textbf{que} institui o Código de Trânsito Brasileiro, nele fica definido \textbf{que} a educação para o trânsito é direito de todos e constitui dever prioritário para os componentes do Sistema Nacional de Trânsito : > Fica instituída \textbf{que} a educação para o trânsito será promovida na pré-escola e nas escolas de 1º, 2º e 3º graus, por meio de planejamento e ações coordenadas entre os órgãos e entidades do Sistema Nacional de Trânsito e de Educação, da União, dos Estados, do Distrito Federal e dos Municípios, nas respectivas áreas de atuação.
A adoção, em todos os níveis de ensino, de \textbf{um} currículo interdisciplinar com conteúdo programático sobre segurança de trânsito ( CTB, 1997 ).
No estado do Ceará esta parceria vem sendo desenvolvida através da cooperação intersetorial entre a Secretaria de Saúde, a Secretaria de Educação e os demais órgãos competentes. e ) Cidadania e Civismo O tema transversal cidadania já aparece nos Parâmetros Curriculares Nacionais - PCNS ( 2000 ), dentro da macroárea Trabalho e Consumo e, deste modo, \textbf{permeia} as diferentes áreas do conhecimento.
Na BNCC ( 2018 ) o tema é adaptado para a realidade mais \textbf{atual} e relacionado à ideia de civismo.
Ele exprime conceitos e valores básicos à democracia, à participação política em defesa do estado de direitos sociais e ao dever individual de cada \textbf{um} em reconhecer a importância de proferir esses valores associados à condição de cidadão cônscio e responsável direto pelo pleno vigor desses preceitos.
Assim, o indivíduo é capaz de discutir direitos, deveres e valores importantes e \textbf{urgentes} para o fortalecimento da democracia na sociedade brasileira \textbf{contemporânea}, respeitando as diferenças de \textbf{uma} sociedade \textbf{que} é múltipla, complexa, de modo inclusivo e tolerante.
Na BNCC, a temática está relacionada diretamente a algumas das competências gerais da educação básica ( \textbf{notadamente} as de número 6, 7, 9 e 10 ), valorizando a diversidade de saberes e as vivências culturais ; promovendo o exercício da cidadania, os direitos humanos, o diálogo e o respeito mútuo com urgência social e abrangência nacional ; além de possibilitar \textbf{uma} aprendizagem significativa, \textbf{que} favoreça o entendimento da realidade social e \textbf{fomente} a participação política e cidadã com base em princípios éticos, democráticos, inclusivos, sustentáveis e solidários.
Desse modo, o tema transversal \textbf{contemporâneo} cidadania e civismo \textbf{implica} não somente o aprender – no sentido de reproduzir o conhecimento socialmente adquirido –, mas oportuniza também apreender significativamente a realidade posta para transformá -la.
Dentre os principais subtemas abordados na macroárea, tem -se : aspectos da vida familiar e social ; educação para o trânsito ; educação em direitos humanos ; direitos da criança e do adolescente ; processo de envelhecimento, respeito e valorização do idoso.
No DCRC, esses subtemas dialogam de modo intradisciplinar quando se interrelacionam cidadania e civismo, através de \textbf{um} dos componentes curriculares ; interdisciplinar, quando se discute os subtemas, de \textbf{um} modo associado, em mais de \textbf{um} componente da mesma área e transdisciplinar, quando \textbf{um} ou mais temas são discutidos para além da individualidade de cada área do conhecimento. f ) Multiculturalidade ( ou Multiculturalismo ), Interculturalidade, Intraculturalidade e Transculturalidade A Multiculturalidade ( ou multiculturalismo ), compreendida como o resultado do encontro ou contato e convivência das diversas culturas de \textbf{uma} sociedade, afigura -se como aporte inicial para pensar \textbf{uma} prática pedagógica adequada às particularidades de cada cidadão e \textbf{que} deve ser refletido nos espaços educacionais, \textbf{exigindo} \textbf{que} seja abordado o conhecimento sobre os demais conceitos relacionados à temática, tais como : interculturalidade, intraculturalidade e transculturalidade, já \textbf{que} pensar somente o multiculturalismo por si só, poderá reforçar a ideia de \textbf{uma} cultura dominante ou \textbf{hegemônica}, sem nenhuma preocupação com o processo de descolonização, pela própria origem do conceito e seus pensadores.
A priori, olhar a escola sob o prisma da multiculturalidade é reconhecer a diversidade cultural \textbf{que} ali se expressa através dos estudantes, dando início ao exercício da cidadania, haja vista o reconhecimento da identidade cultural pelos membros de \textbf{uma} coletividade ser a base para a efetivação da cidadania. \textbf{Uma} perspectiva educacional \textbf{amparada} sob pressupostos multiculturais pensa os aspectos da diversidade cultural em qualquer sociedade, seu pluralismo cultural ou diversidade étnica e cultural, \textbf{que} no Brasil tem os índios, negros e europeus como elementos \textbf{constitutivos}.
As competências gerais \textbf{que} subsidiam a construção da BNCC e o desenvolvimento de qualquer currículo escolar devem ser diretamente dependentes da \textbf{ótica} dos conceitos sugeridos neste subitem, pois responsabilidade e cidadania, empatia e cooperação, autoconhecimento e autocuidado, bem como o próprio conhecimento cultural dos diversos povos \textbf{que} constituem o Brasil e o Ceará serão base para \textbf{que} essa visão se manifeste. \textbf{Uma} prática pedagógica \textbf{que} não observa esses aspectos caminha para o insucesso do processo educacional dos estudantes e o enfraquecimento da cidadania, haja vista \textbf{que} despotencializa a formação de subjetividades \textbf{que} se autorreconhece com identidades singulares e precisam ser respeitadas em suas diferenças.
Multiculturalidade e interculturalidade são conceitos \textbf{que} se apresentam como \textbf{um} ponto fundamental : \textbf{que} a sociedade não é homogênea e sim plural. \textbf{Contudo}, há \textbf{uma} diferença primordial, a multiculturalidade manifesta as diferenças como \textbf{um} fato objetivo e a interculturalidade \textbf{inspira} \textbf{que} estas diferenças dialoguem em plena tradução das culturas.
No presente, entretanto, a palavra interculturalidade deixa de ser \textbf{um} \textbf{mero} diálogo entre as culturas e passa a definir projetos ético-políticos voltados à construção de relações sociais horizontais entre diferentes grupos humanos, sem \textbf{que} seja \textbf{preciso} abdicar da diferença.
Pelo contrário, a marcação da diferença é \textbf{um} dos traços da \textbf{contemporaneidade}.
Se o processo de colonização investiu pesadamente na homogeneização cultural, o \textbf{atual} momento é de valorização das diferenças, com \textbf{um} consequente retorno ao interior das culturas, num processo de autorreconhecimento e de posterior saída para o diálogo \textbf{baseado} na alteridade.(MARTINS, 2021. p.
XIII ) A interculturalidade é \textbf{um} conjunto de propostas de convivência democrática entre diferentes culturas, ou seja, se refere a diversidade cultural. \textbf{Vivemos} isso no nosso dia a dia, nas grandes cidades.
Com isso, vamos encontrar pessoas de variadas origens étnicas, línguas diferentes e tradições culturais bem diversificadas, buscando a integração entre elas sem anular sua diversidade, ao contrário, promovendo a integração e \textbf{uma} aposta na aceitação e na normalização das diferenças sociais.
Outros conceitos importantes são a intraculturalidade e a transculturalidade.
Gervás ( 2011. p. 15 ) entende a intraculturalidade como o processo de autorreconhecimento da identidade cultural e étnica pelo qual o indivíduo passa : > Ou seja, olhando culturalmente para dentro de sua própria pessoa e de sua própria cultura, tentando conhecer -se e valorizar -se socialmente e culturalmente a nós mesmos, através da complexidade e diferença interna do próprio grupo social. \textbf{Uma} vez \textbf{que} este objetivo é alcançado, nós podemos considerar, então, abrangendo aspectos interculturais ou multiculturais.
E a transculturalidade é justamente a finalização do processo de interação entre grupos culturais, gerando alguma transformação cultural sobre \textbf{um} ou mais grupos envolvidos.
Nesse sentido, através do multiculturalismo pode -se promover a interculturalidade, em busca da intraculturalidade e assim fazer surgir a transculturalidade ( Martins, 2015 ), construindo modelos de intervenções culturais decolonizantes, deixando no passado a cópia de formatos eurocêntricos e etnocêntricos e percebendo \textbf{que} a multiculturalidade, a interculturalidade, a intraculturalidade e a transculturalidade devem fazer parte do nosso cotidiano, o \textbf{que} viabiliza processos educativos em sintonia com o presente e fundamentais aos estudantes, professores e toda a sociedade. g ) Ciência e Tecnologia O desenvolvimento da ciência e da tecnologia acarreta transformações na sociedade, refletindo em mudanças nos níveis econômico, político e social.
Comumente, a ciência e a tecnologia são \textbf{vistas} como motores relacionados aos conceitos de progresso e desenvolvimento, oriundos da visão ocidental. \textbf{Contudo}, \textbf{uma} visão positivista de ciência pode fortalecer o mito da cientificidade, confiança excessiva na ciência e na tecnologia como campos de conhecimento isolados e \textbf{que} seriam as únicas com respostas adequadas à resolução dos diversos problemas da sociedade.
Os interesses sociais, políticos, econômicos, de segurança e defesa nacional \textbf{implicam} eventuais riscos, porquanto o desenvolvimento científico-tecnológico e seus produtos não são independentes de seus interesses.
Bazzo ( 1998, p. 142 ) destaca \textbf{que} : > É inegável a contribuição \textbf{que} a ciência e a tecnologia trouxeram nos últimos anos.
Porém, apesar desta constatação, não podemos confiar excessivamente nelas, tornando–nos cegos pelo conforto \textbf{que} nos proporcionam cotidianamente seus aparatos e dispositivos técnicos.
Isso pode resultar perigoso porque, nesta anestesia \textbf{que} o deslumbramento da modernidade tecnológica nos oferece, podemos nos esquecer \textbf{que} a ciência e a tecnologia incorporam questões sociais, éticas e políticas.
Torna–se cada vez mais necessário \textbf{que} a população possa ter condições de acessar, avaliar e participar das decisões e informações sobre o desenvolvimento científico–tecnológico \textbf{que} venham a atingir o meio onde \textbf{vivem}.
É necessário \textbf{que} a sociedade, em geral, comece a questionar sobre os impactos da evolução e aplicação da ciência e tecnologia sobre seu entorno e \textbf{consiga} perceber \textbf{que}, muitas vezes, certas atitudes não atendem à maioria, mas, sim, aos interesses dominantes.
A esse respeito, Bazzo ( 1998, p. 34 ) comenta :"o cidadão merece aprender a ler e entender muito mais do \textbf{que} conceitos estanques – a ciência e a tecnologia, com suas implicações e conseqüências, para poder ser elemento participante nas decisões de ordem política e social \textbf{que} \textbf{influenciarão} o seu futuro e o dos seus filhos".
O Enfoque Ciência-Tecnologia-Sociedade ( CTS ) propõe práticas e metodologias no ensino das ciências com base no contexto educativo, o \textbf{que} ressalta a necessidade de colocar ciência e tecnologia em novas concepções vinculadas ao contexto social.
De acordo com Medina e Sanmartín ( 1990 ), quando se pretende incluir o enfoque CTS no contexto educacional, alguns objetivos precisam ser buscados : - Questionar as formas herdadas de estudar e atuar sobre a natureza, as quais devem ser constantemente refletidas.
Sua legitimação deve ser feita por meio do sistema educativo, pois só assim é possível contextualizar permanentemente os conhecimentos em função das necessidades da sociedade. - Questionar a distinção convencional entre conhecimento \textbf{teórico} e conhecimento prático – assim como sua distribuição social entre ' os \textbf{que} pensam ' e ' os \textbf{que} executam ' – \textbf{que} reflete, por sua vez, \textbf{um} sistema educativo dúbio, \textbf{que} diferencia a educação geral da vocacional. - Combater a segmentação do conhecimento, em todos os níveis de educação. - Promover \textbf{uma} autêntica \textbf{democratização} do conhecimento científico e tecnológico, de modo \textbf{que} ela não só se difunda, mas \textbf{que} se integre na atividade produtiva das comunidades de maneira crítica.
Dessa forma, a importância de discutir com os alunos os avanços da ciência e tecnologia, suas causas, consequências, os interesses econômicos e políticos, de forma contextualizada, está no fato de \textbf{que} devemos conceber a ciência como fruto da criação humana, intimamente ligada à nossa evolução, \textbf{permeada} pela ação reflexiva de quem sofre/age as diversas \textbf{crises} \textbf{inerentes} a esse processo de desenvolvimento ( PINHEIRO ; SILVEIRA ; BAZZO, 2007 )

% matched lemmas: abordagem, amparar, argumentar, atual, basear, cientista, configurar, conhecido, conseguinte, conseguir, constitutivo, contemporaneidade, contemporâneo, contudo, crise, democratização, despertar, direção, enquanto, exigir, explicitar, fomentar, globalizado, hegemônico, homem, humanidade, implicar, inerente, influenciar, inspirar, justo, mero, metodológico, notadamente, partilha, pedagogia, permear, pontuar, posicionamento, preciso, que, raça, restringir, sujeito, século, sólido, teórico, transdisciplinaridade, um, urgente, ver, viver, ótica
\end{textsample}
